\documentclass[10pt]{article}

% ---------- Document Identity (update these only) ----------
\newcommand{\DocStage}{}
\newcommand{\DocTitle}{Your Road to Tier One SOF}
\newcommand{\ProgramName}{182-Week Civilian-to-Selection Readiness Pipeline}
\newcommand{\ProgramSubtitle}{}

% ---------- Page ----------
\usepackage[legalpaper,left=0.5in,right=0.5in,top=0.6in,bottom=0.45in]{geometry}

% ---------- Typography ----------
\usepackage[T1]{fontenc}
\usepackage[utf8]{inputenc}
\usepackage{libertinus}
\usepackage{microtype}
\usepackage{parskip}
\usepackage{ragged2e}
\usepackage{textcomp}
\AtBeginDocument{\color{black}}
\AtBeginDocument{\pagecolor{white}}
\AtBeginDocument{\justifying}

% ---------- Landscape pages ----------
\usepackage{pdflscape}

% ---------- Color ----------
\usepackage[table]{xcolor}
\definecolor{StageBlue}{HTML}{1F4E79}
\definecolor{StageBlueDark}{HTML}{163A5A}
\definecolor{LightGray}{HTML}{F2F4F7}
\definecolor{MidGray}{HTML}{D9DEE5}
\definecolor{RuleGray}{HTML}{C7CED8}
\definecolor{HeaderInk}{HTML}{000000}
\definecolor{Ink}{HTML}{000000}

% ---------- Utility ----------
\usepackage{enumitem}
\sloppy
\setlength{\emergencystretch}{2em}

% ---------- Boxes / UI ----------
\usepackage[most]{tcolorbox}

\tcbset{
  charterTitle/.style={
    enhanced,
    colback=StageBlue!4,
    colframe=StageBlue,
    boxrule=1.25pt,
    arc=3pt,
    left=16pt,right=16pt,top=14pt,bottom=14pt,
    borderline north={3.2pt}{0pt}{StageBlue},
    borderline west={4pt}{0pt}{StageBlue},
    borderline south={1.25pt}{0pt}{StageBlue},
    drop shadow={black!25!white},
  },
  charterRule/.style={
    enhanced,
    colback=LightGray,
    colframe=RuleGray,
    boxrule=0.85pt,
    arc=3pt,
    left=12pt,right=12pt,top=10pt,bottom=10pt,
    borderline west={3pt}{0pt}{StageBlue},
  },
  charterBadge/.style={
    enhanced,
    colback=StageBlue,
    coltext=white,
    colframe=StageBlueDark,
    boxrule=0.6pt,
    arc=2pt,
    left=6pt,right=6pt,top=3pt,bottom=3pt,
  }
}
\newtcbox{\Badge}{nobeforeafter,charterBadge}

% ---------- Lists ----------
\setlist[itemize]{leftmargin=1.5em, itemsep=0.25em, topsep=0.25em}
\setlist[enumerate]{leftmargin=1.8em, itemsep=0.25em, topsep=0.25em}

% ---------- TOC ----------
\usepackage{tocloft}
\setcounter{tocdepth}{3}
\setlength{\cftbeforesecskip}{3pt}
\setlength{\cftbeforesubsecskip}{2pt}
\setlength{\cftbeforesubsubsecskip}{0pt}
\renewcommand{\cftsecfont}{\bfseries}
\renewcommand{\cftsecpagefont}{\bfseries}
\renewcommand{\cftsubsecfont}{\normalfont}
\renewcommand{\cftsubsecpagefont}{\normalfont}
\renewcommand{\cftsubsubsecfont}{\normalfont}
\renewcommand{\cftsubsubsecpagefont}{\normalfont}
\setlength{\cftsecindent}{0pt}
\setlength{\cftsubsecindent}{1.25em}
\setlength{\cftsubsubsecindent}{2.8em}
\setlength{\cftsecnumwidth}{2.1em}
\setlength{\cftsubsecnumwidth}{2.8em}
\setlength{\cftsubsubsecnumwidth}{3.6em}

% Remove the built-in TOC heading so only the custom heading is shown
\makeatletter
\renewcommand{\tableofcontents}{\@starttoc{toc}}
\makeatother

% ---------- Header/Footer: page number only ----------
\usepackage{fancyhdr}
\pagestyle{fancy}
\fancyhf{}
\renewcommand{\headrulewidth}{0pt}
\renewcommand{\footrulewidth}{0pt}
\fancyfoot[C]{\thepage}

% ---------- Section formatting ----------
\usepackage{titlesec}

\titleformat{\section}
  {\Large\bfseries\color{Ink}}
  {\thesection}{0.6em}{}
  [\vspace{0.25em}\textcolor{StageBlue}{\rule{\linewidth}{0.8pt}}]

\titleformat{\subsection}
  {\large\bfseries\color{Ink}}
  {\thesubsection}{0.6em}{}

\titleformat{\subsubsection}
  {\normalsize\bfseries\color{Ink}}
  {\thesubsubsection}{0.6em}{}

\titlespacing*{\section}{0pt}{0.95em}{0.35em}
\titlespacing*{\subsection}{0pt}{0.65em}{0.25em}
\titlespacing*{\subsubsection}{0pt}{0.55em}{0.2em}

% ---------- Table engine ----------
\usepackage{tabularray}
\UseTblrLibrary{booktabs}

% ---------- tabularray: GLOBAL table treatment ----------
\SetTblrInner{
  cells = {fg=black, bg=white, valign=t},
}

% ---------- Header helpers ----------
\newcommand{\Hdr}[1]{\shortstack[c]{\strut #1\strut}}

% ---------- Lines ----------
\newcommand{\TitleLine}{\textcolor{StageBlue}{\rule{0.98\linewidth}{0.95pt}}}
\newcommand{\SubTitleLine}{\textcolor{RuleGray}{\rule{0.98\linewidth}{0.65pt}}}

% ---------- Continued on next page ----------
\newcommand{\ContinuedOnNextPageInline}{%
  \par
  \vfill
  \noindent\textcolor{RuleGray}{\rule{\linewidth}{1pt}}\vspace{-2pt}
  \noindent\hfill\textit{Continued on next page}\par\vspace{-10pt}
  \noindent\textcolor{RuleGray}{\rule{\linewidth}{1pt}}
  \clearpage
}

% ---------- PDF hyperlinks ----------
\usepackage[hidelinks]{hyperref}
\hypersetup{
  colorlinks=false,
  pdfborder={0 0 0},
  linktoc=all,
  pdftitle={\DocTitle},
  pdfauthor={\ProgramName}
}

% =========================================================
\begin{document}

\section*{7.16 — Operational Self-Management Systems — Journaling, Planning, Mental Conditioning}
\phantomsection
\addcontentsline{toc}{section}{7.16 — Operational Self-Management Systems — Journaling, Planning, Mental Conditioning}

\subsection*{7.16.1 Purpose \& Scope}
\phantomsection
\addcontentsline{toc}{subsection}{7.16.1 Purpose \& Scope}

7.16 provides the operational systems that make adherence and auditability possible: planning, journaling, review cadence, and non-clinical mental skills routines that keep execution safe and consistent. It supports the required datasets for daily/session/test logs and enforces the “no inference” posture for missing data.

\subsection*{7.16.2 Governance and Safety Boundaries}
\phantomsection
\addcontentsline{toc}{subsection}{7.16.2 Governance and Safety Boundaries}

\begin{itemize}
\item Non-clinical boundary: this is performance and adherence infrastructure, not mental health treatment.
\item Evidence integrity boundary: missing or reconstructed logs degrade admissibility for gates; record \textbf{MISSING}/\textbf{UNKNOWN} rather than inferring.
\item Stop posture boundary: mental \textbf{STOP} \textbf{NOW} triggers must be treated as safety events, logged, and escalated for professional support at a high level.
\item Privacy boundary: record only the minimum necessary; protect the log store; avoid oversharing sensitive content.
\end{itemize}

\subsection*{7.16.3 Tier Doctrine — Applied to Self-Management}
\phantomsection
\addcontentsline{toc}{subsection}{7.16.3 Tier Doctrine — Applied to Self-Management}

\begin{itemize}
\item Tier 0: start-from-zero capture options; minimal friction; offline-first by default; sufficient to begin compliance.
\item Tier 1: minimum deployable system-of-record with offline fallback, repeatable templates, and weekly audit routine.
\item Tier 2: expanded structure: dashboards/rollups, more robust staging scripts, and stronger reslot discipline tooling as consequence increases.
\item Tier 3: redundancy and high reliability: duplicate capture pathways, travel kits, and structured decision-window controls that reduce variance.
\end{itemize}

\subsection*{7.16.4 Minimum Dataset Reminder — Stage 2.E field list only}
\phantomsection
\addcontentsline{toc}{subsection}{7.16.4 Minimum Dataset Reminder — Stage 2.E field list only}

Daily minimum fields:

\begin{itemize}
\item ISO timestamp
\item Phase/Week
\item Sleep hours
\item Sleep quality 0–10
\item Soreness 0–10
\item Fatigue 0–10
\item Pain 0–10 plus location
\item Steps count
\item Notes
\end{itemize}

Session minimum fields:

\begin{itemize}
\item ISO timestamp
\item Phase/Week
\item Session \#
\item Session Validity Tag (\textbf{VALID} / \textbf{MODIFIED} / \textbf{INVALID})
\item \textbf{SESSION-RC} (only if \textbf{MODIFIED} or \textbf{INVALID})
\item Cardio Min Standard Met (Y/N)
\item Cardio modality
\item Cardio duration minutes
\item Cardio intensity (RPE plus talk test cue)
\item Main work domains
\item Session RPE 0–10
\item Pain flags
\item Foot flags
\item Notes
\end{itemize}

Test-event minimum fields:

\begin{itemize}
\item ISO timestamp
\item Phase/Week
\item Protocol Card C\#
\item Linked Metric \textbf{ID}(s)
\item Domain tag
\item Environment unit \textbf{ID}
\item Tool standard (timer/scale/tape)
\item Warm-up completed (Y/N)
\item Test Validity Tag (\textbf{VALID} / \textbf{INVALID})
\item \textbf{TEST-RC} (only if \textbf{INVALID})
\item Result
\item Evidence Ref
\item Replication status
\item Conditions notes
\item Stop/Abort (Y/N plus reason)
\end{itemize}

\subsection*{7.16.5 Selection Criteria \& Fit Checks}
\phantomsection
\addcontentsline{toc}{subsection}{7.16.5 Selection Criteria \& Fit Checks}

\begin{itemize}
\item Low friction: the system must be usable daily without “catch-up” burden.
\item Single source of truth: one canonical repository; avoid parallel logs that diverge.
\item Offline resilience: paper or offline file must exist for device failures.
\item Privacy posture: protect sensitive entries; record only what is required.
\item Fit checks:
  \begin{itemize}
  \item daily log completion time is short and repeatable,
  \item session logs can be completed same-day,
  \item test-event logs capture evidence refs and environment IDs without reconstruction.
  \end{itemize}
\end{itemize}

\subsection*{7.16.6 Setup, Safety, and Anchoring}
\phantomsection
\addcontentsline{toc}{subsection}{7.16.6 Setup, Safety, and Anchoring}

\begin{itemize}
\item Setup posture:
  \begin{itemize}
  \item designate one canonical system-of-record location and one offline fallback location.
  \item pre-build templates for daily/session/test logs so no new schema is invented mid-execution.
  \end{itemize}
\item Safety anchoring:
  \begin{itemize}
  \item stop posture is logged as a safety event; do not hide it to protect pride or narrative.
  \end{itemize}
\item Protected window anchoring:
  \begin{itemize}
  \item use pre-window checklists to freeze variables (routes/tools/shoes/caffeine changes) per governance.
  \end{itemize}
\item Non-clinical mental skills anchoring:
  \begin{itemize}
  \item short, repeatable routines (breathing, attention control, planning) are permitted; avoid clinical framing.
  \end{itemize}
\end{itemize}

\subsection*{7.16.7 Maintenance, Replacement, and Storage}
\phantomsection
\addcontentsline{toc}{subsection}{7.16.7 Maintenance, Replacement, and Storage}

\begin{itemize}
\item Weekly maintenance:
  \begin{itemize}
  \item review missing fields; mark \textbf{MISSING}/\textbf{UNKNOWN} explicitly; do not infer.
  \item archive weekly summaries.
  \end{itemize}
\item Backup posture:
  \begin{itemize}
  \item periodic export/backup of digital logs to prevent loss.
  \item store paper logs dry and organized.
  \end{itemize}
\item Template maintenance:
  \begin{itemize}
  \item update only via patch discipline; avoid “silent schema drift.”
  \end{itemize}
\end{itemize}

\subsection*{7.16.8 Substitution Rules — When Tools Fail}
\phantomsection
\addcontentsline{toc}{subsection}{7.16.8 Substitution Rules — When Tools Fail}

\begin{itemize}
\item If the primary digital system fails:
  \begin{itemize}
  \item switch to offline fallback immediately; capture the same required fields; later transcribe without memory inference.
  \end{itemize}
\item If a weekly review is missed:
  \begin{itemize}
  \item perform it at the next opportunity; do not backfill metrics—use available evidence only.
  \end{itemize}
\item If evidence refs are missing:
  \begin{itemize}
  \item record as \textbf{MISSING}; do not claim gate credit where evidence is required; reslot tests if needed.
  \end{itemize}
\item Validity linkage:
  \begin{itemize}
  \item deviations that affect admissibility require correct validity tagging and reason-code posture per Stage 1.F; do not invent new reason codes.
  \end{itemize}
\end{itemize}

\subsection*{Patch Record Template — aligned to Stage 1.F}
\phantomsection
\addcontentsline{toc}{subsection}{Patch Record Template — aligned to Stage 1.F}

\begin{landscape}

\begingroup
\scriptsize
\setlength{\tabcolsep}{2pt}
\renewcommand{\arraystretch}{1.12}

\begin{longtblr}{
  width=\linewidth,
  colspec = {
    Q[c,wd=0.90in]
    Q[c,wd=0.95in]
    X[c,2.20]
    X[c,1.40]
    X[c,1.75]
    X[c,1.75]
    Q[c,wd=0.95in]
    Q[c,wd=1.05in]
  },
  rowhead=1,
  hlines,
  vlines,
  cells = {hsep=2pt, vsep=6pt, halign=c, valign=m},
  row{1} = {bg=StageBlue!15, fg=black, font=\bfseries, valign=m},
  row{odd} = {bg=white},
  row{even} = {bg=LightGray},
}
\Hdr{Patch \textbf{ID}} &
\Hdr{Date} &
\Hdr{Target (Stage/Section)} &
\Hdr{Change} &
\Hdr{Rationale} &
\Hdr{Validation Check} &
\Hdr{Owner} &
\Hdr{Impact Level} \\
\end{longtblr}

\endgroup

\end{landscape}

\subsection*{Weekly Audit Checklist Template}
\phantomsection
\addcontentsline{toc}{subsection}{Weekly Audit Checklist Template}

\begin{itemize}
\item Session validity distribution reviewed (\textbf{VALID} / \textbf{MODIFIED} / \textbf{INVALID}).
\item Test validity rate reviewed (\textbf{VALID} / \textbf{INVALID}).
\item Tool consistency logged (tool changes explicitly recorded).
\item Route/environment consistency logged (Route \textbf{ID} / Environment Unit \textbf{ID} stability where applicable).
\item Risk flags reviewed (pain trends, foot flags, sleep/fatigue trends).
\item Upcoming Deload+Test readiness pointers reviewed (no novelty posture; staging complete).
\item Reslot log reviewed for missed/invalid tests; no make-up stacking; next protected window identified.
\item Patch log reviewed: any changes recorded; no silent edits.
\end{itemize}

\subsection*{7.16.9 Phase Integration Map — Required}
\phantomsection
\addcontentsline{toc}{subsection}{7.16.9 Phase Integration Map — Required}

\begin{landscape}

\begingroup
\scriptsize
\setlength{\tabcolsep}{2pt}
\renewcommand{\arraystretch}{1.12}

\begin{longtblr}{
  width=\linewidth,
  colspec = {
    Q[c,wd=1.05in]
    Q[c,wd=1.55in]
    X[c,3.40]
    X[c,3.85]
  },
  rowhead=1,
  hlines,
  vlines,
  cells = {hsep=2pt, vsep=6pt, halign=c, valign=m},
  row{1} = {bg=StageBlue!15, fg=black, font=\bfseries, valign=m},
  row{odd} = {bg=white},
  row{even} = {bg=LightGray},
}
\Hdr{Phase} &
\Hdr{Required Tier (from Stage 6)} &
\Hdr{What Changes / Becomes Mandatory} &
\Hdr{Notes} \\
Phase 00 &
T1 (E) &
Tier 1 required: single system-of-record + offline fallback; daily and session logs executed same-day &
Phase 00 (E): minimum viable evidence system is required for admissibility posture. \\
Phase 01 &
T1 (End) &
Maintain Tier 1; add routine weekly audit cadence as volume rises &
Phase 01 (End): weekly audit becomes operationally necessary; templates remain stable. \\
Phase 02 &
T1 (End) &
Maintain Tier 1; stricter “no novelty near protected windows” checklists as test consequence increases &
Phase 02 (End): pre-window lock scripts become higher ROI. \\
Phase 03 &
T1 (End) &
Maintain Tier 1; ruck/systems complexity increases need for staging scripts and reslot discipline &
Phase 03 (End): increase check-in script specificity (pain/foot flags). \\
Phase 04 &
T1 (End) &
Maintain Tier 1; multi-domain logistics increases audit discipline &
Phase 04 (End): weekly audit becomes non-negotiable for drift detection. \\
Phase 05 &
T1 (End) &
Maintain Tier 1; longer efforts increase value of dashboard rollups (optional until Stage 6 elevates tier) &
Phase 05 (End): add simple rollups if missed sessions or drift appears. \\
Phase 06 &
T2 (End) &
Tier 2 becomes required: dashboard rollups + reslot decision log + scripted daily check-in &
Phase 06 (End): consequence and density justify stronger tooling; prevent make-up stacking. \\
Phase 07 &
T2 (End) &
Maintain Tier 2; increase audit cadence tightness near protected windows &
Phase 07 (End): weekly audit executed with higher discipline; reslot log actively used. \\
Phase 08 &
T2 (End) &
Maintain Tier 2; increase check-in emphasis for durability/foot interface &
Phase 08 (End): pain/foot flags drive downshift/reslot decisions; keep reporting honest. \\
Phase 09 &
T2 (End) &
Maintain Tier 2; increase check-in specificity for high-volume fatigue &
Phase 09 (End): dashboard rollups help prevent silent overload. \\
Phase 10 &
T2 (End) &
Maintain Tier 2; multi-domain gates increase value of pre-window lock scripts &
Phase 10 (End): expand lock script to include tool/route/footwear and stimulant stability. \\
Phase 11 &
T2 (End) &
Maintain Tier 2; recovery sensitivity increases audit criticality &
Phase 11 (End): audit prevents invalid gate interpretations due to fatigue contamination. \\
Phase 12 &
T3 (End) &
Tier 3 becomes required: redundant capture pathway + travel kit + decision-window lock script &
Phase 12 (End): high-consequence windows justify redundancy to prevent data loss and variance. \\
Phase 13 &
T3 (End) &
Maintain Tier 3; consolidation emphasizes low variance and repeatable evidence &
Phase 13 (End): lock scripts and redundancy prevent last-minute drift. \\
\end{longtblr}

\endgroup

\end{landscape}

7.16.9 Phase Integration Map Notes:

\begin{itemize}
\item If any phase artifact demands Tier 2/3 self-management tools earlier than Stage 6 timing, Requires Stage 8 review.
\item If any artifact attempts to change Stage 2.E required fields or Stage 1.F validity tags via “convenience,” Requires Stage 8 review.
\end{itemize}

\subsection*{7.16.10 Risks, Contraindications, and Red Flags}
\phantomsection
\addcontentsline{toc}{subsection}{7.16.10 Risks, Contraindications, and Red Flags}

\begin{itemize}
\item Evidence integrity risks:
  \begin{itemize}
  \item missing logs, reconstructed entries, and multi-source log drift can invalidate gate decisions.
  \item control: single source-of-truth + offline fallback; \textbf{MISSING}/\textbf{UNKNOWN} posture; weekly audit.
  \end{itemize}
\item Behavioral drift risks:
  \begin{itemize}
  \item missed sessions “made up” by stacking violates reslot discipline and increases injury risk.
  \item control: reslot log and weekly audit; enforce no-stacking posture.
  \end{itemize}
\item Privacy risks:
  \begin{itemize}
  \item oversharing personal content can reduce compliance; keep records minimal and operational.
  \end{itemize}
\item Mental \textbf{STOP} \textbf{NOW} triggers (explicit):
  \begin{itemize}
  \item acute panic or dissociation → \textbf{STOP} \textbf{NOW}; stop exposure; log event; seek professional support.
  \item suicidal ideation or self-harm statements → \textbf{STOP} \textbf{NOW}; log; seek professional support.
  \item inability to safely self-regulate in-session → \textbf{STOP} \textbf{NOW}; stop exposure; log; seek professional support.
  \end{itemize}
\item Red-flag posture:
  \begin{itemize}
  \item do not test that day when \textbf{STOP} \textbf{NOW} triggers occur.
  \item do not attempt to “work through” severe mental distress as training.
  \end{itemize}
\item Requires Stage 8 review triggers:
  \begin{itemize}
  \item any downstream artifact that treats mental health events as “normal discomfort” rather than \textbf{STOP} \textbf{NOW}.
  \item any timing mismatch where a phase requires capabilities earlier than Stage 6 Tier timing.
  \end{itemize}
\end{itemize}

\subsection*{7.16.11 Compliance Checklist — Pass/Fail}
\phantomsection
\addcontentsline{toc}{subsection}{7.16.11 Compliance Checklist — Pass/Fail}

\begin{itemize}
\item \textbf{PASS}/\textbf{FAIL}: Tier 1 includes offline-capable fallback and is required by Phase 00 (E) per Stage 6 timing.
\item \textbf{PASS}/\textbf{FAIL}: Minimum Dataset Reminder block lists daily/session/test-event fields without altering schema.
\item \textbf{PASS}/\textbf{FAIL}: Patch Record Template matches Stage 1.F posture and includes required fields.
\item \textbf{PASS}/\textbf{FAIL}: Weekly Audit Checklist Template aligns to Stage 1.F and Stage 3.E posture and introduces no new numeric rules.
\item \textbf{PASS}/\textbf{FAIL}: Phase Integration Map aligns to Stage 6 timing and calls out later-phase increases in check-in scripts, audit cadence, and dashboard rollups.
\item \textbf{PASS}/\textbf{FAIL}: Mental \textbf{STOP} \textbf{NOW} triggers are explicit and include stop, log, and seek professional support posture.
\end{itemize}

\end{document}